\section{Discussion}\label{sec:discussion}

\paragraph{Additional Related Work.} \cite{zhang2024predictable} uses proxy benchmarks with small proxy models to predict pre-training performance of larger target models. However, their approach suffers from several limitations. First, it is computationally expensive, requiring global search over 42 benchmarks and 34 models. Second, the required size of the search space (number of benchmarks and models) lacks principled justification. Third, using fundamentally different test distributions introduces suboptimal noise, as we empirically demonstrate in Appendix \ref{app:proxy_task}. 

\cite{xie2023doremi} and \cite{liu2025regmix} leverage small proxy models (1 to 280M parameters) to optimize pre-training data \textit{mixture weights} for larger target models ($\sim$7B), operating under the assumption that small-scale performance meaningfully correlates with large-scale performance. Both works rely on standard NLL or perplexity metrics, which we demonstrate to be suboptimal in \textbf{\S\ \ref{sec:limit}}. Furthermore, their experiments are predominantly on non-emergent benchmarks (e.g., TriviaQA and HellaSwag), whereas our work specifically targets reasoning benchmarks that exhibit emergent behavior. Lastly, their approaches are specialized for optimizing pre-training data \textit{mixtures}, whereas our method addresses the more general problem of predicting large model performance with small proxy models, enabling \textit{broader} applications.

Moreover, readers can view \cite{10651521}, \cite{10.1145/3706118}, \cite{kipnis2025metabench}, and \cite{pacchiardi-etal-2025-predictaboard} for some past works on LLM performance prediction, and evaluation. We acknowledge that alternative branches of the literature have studied predicting LLM performance.

\paragraph{Minimal Compute Overhead.} While \tsc{rBridge} introduces some additional costs, these are minor. First, is generating the gold reasoning trace $R^\phi$ via the frontier model. However, this represents a small one-time cost of under \$10 per benchmark. To help the research community we plan to open-sourced our dataset, allowing future researchers to skip this cost for benchmarks we examine. Second, the automatic weighting mechanism (Fig. \ref{fig:rbridge}) incurs only a few seconds of CPU runtime per benchmark. Both overheads are negligible compared to the computational savings our method provides.
% \paragraph{Technical Takeaway: Pre-training Evaluation Objective and Task Alignment.} Our main technical finding is that we should attempt to find an output $Y^*$ that is more ID while ensuring that it is task aligned. We first demonstrate that we can increase level of ID by using a frontier model's reasoning trace $R^\phi$ as $Y^*$ (\S\ \ref{sec:easy}, Fig. \ref{fig:idood}). Our empirical finding here is that we should minimize any formatting that is unnatural to the pre-training distribution, e.g., "Final Answer: $\cdots$". Furthermore, $R^\phi$ is well aligned with the task, as it is the intermediary process that leads to the downstream final answer (Fig. \ref{fig:scb_clean}, Tab. \ref{tab:scb}). Then, we further improve task alignment by automatically weighting each token by its level of task alignment during NLL evaluation (\S\ \ref{sec:target_align}, Fig. \ref{fig:diagram}). 

\paragraph{\tsc{rBridge} is a Better Predictor of Large Scale Pre-training Performance on Reasoning.} While we would expect the \textit{same} metric to yield the greatest relationship and predictive power, our paper demonstrates that this is not the case. As visualized in Fig. \ref{fig:target_metric} even as we scale the proxy model by 7$\times$, 13$\times$, it fails to achieve the performance levels of \tsc{rBridge}. Our work empirically demonstrates that improved alignment with the \textbf{(1)} pre-training evaluation objective, and the \textbf{(2)} task is key to successfully leveraging small models to proxy large model's reasoning performance. Our ablative studies clearly indicate that each component \textbf{(1, 2)} of \tsc{rBridge} is valuable 
(Tab. \ref{tab:scb}, Fig. \ref{fig:ablation}).


\paragraph{Significant Compute and Economic Cost Reduction.} We observe definitive computational cost saving gains through \tsc{rBridge}. In our first experiment \textbf{(i)}, we demonstrate that through \tsc{rBridge} proxy models that are 324.3$\times$ (proxy models $\frac{3.7}{1200}$ the size of target) smaller can be strong proxies for dataset ranking. From a computational perspective, \tsc{rBridge} achieve at least 100.2$\times$ compute savings against the best baseline in achieving the same ranking performance. In our second experiment \textbf{(ii)}, we demonstrate that 13$\times$ and 32$\times$ smaller proxy models can be effectively used as proxies. We believe that \tsc{rBridge}'s utility extrapolates to much greater proxy-target model scale differences. It is also important to note that, by extension, \tsc{rBridge} has the potential to significantly reduce environmental footprint. This has meaningful societal impact as foundation model development has consequential environmental costs \citep{winsta2025hidden,morrison2025holistically}.

% First, \tsc{rBridge} is shown to be state-of-the-art and most reliable across numerous benchmarks and 3$\times$ transfer settings (Tab. \ref{tab:main_tab}, Fig. \ref{fig:target_metric}). Additionally, dramatically improves the pre-training compute needed to search for optimal datasets (Fig. \ref{fig:pareto}, \ref{fig:flops}). 

\paragraph{Enabling Zero-shot Functional Relationship Transfer Across Datasets.} To our knowledge, we are the first to show that the proxy-target function fitted on one pre-training dataset (as shown in experiment \textbf{(ii)}) can be successfully transferred to an alternative one, with no additional fitting. We observe that the improved relationship fitting of \tsc{rBridge} in experiment \textbf{(ii)}'s Tab. \ref{tab:main_tab} extends to improved prediction (MAE) and ranking performance in Tab. \ref{tab:function}. This saves $\frac{m}{n}\times$ compute as any additional pre-training dataset's performance can be reasonably approximated.

\paragraph{Practical Potential Application: Two-stage Dataset Optimization.} 

Practitioners optimizing pre-training data navigate a vast candidate space of data sources and mixtures \citep{xie2023doremi, han2025trillion, liu2025regmix}. The experiments in \textbf{\S\ \ref{sec:exp}} provide a practically applicable framework for such dataset optimization that can be derived as a strong future work. \textbf{[Stage 1]} Given $N$ candidate datasets, following the setting of experiment \textbf{(i)} offers a cost-effective way to filter out $N - k$ poor datasets at e.g. $<$100M proxy scale, reducing the candidate space to $k$. With high decision accuracy of $\sim$80\% (Fig. \ref{fig:decision_acc}), we can expect bad datasets to have been filtered out correctly. Using relatively easier reasoning benchmarks (e.g., ARC-C) suffices for filtering out \textit{weak} datasets. This approach works because (1) strong candidates must also perform well on these tasks, and (2) their performance still correlates with more difficult ones (Fig. \ref{fig:proxytask}), even if they are not strong enough to serve as accurate predictors for ranking the downstream reasoning task. \textbf{[Stage 2]} For the remaining $k$ datasets, we train larger 1B-scale proxies following \textbf{(ii)} and \textbf{(iii)} to accurately rank them by predicting performance at the target scale (e.g., 32B). Although one training run at the target scale $m$ is needed, it is rare for practitioners to work without a baseline target model to improve upon. 

Denoting the required proxy and target compute as $C^\text{p}_1$ and $C^\text{p}_2$ for proxies and $C^\text{t}$ for target, it can be assumed that these costs $C$ are directly proportional to model size \citep{kaplan2020scaling}. The required total compute under this framework is $NC^\text{p}_1 + kC^\text{p}_2 + C^\text{t}$ compared to the naïve approach cost $NC^\text{t}$. Given that the cost reduction factor is $\frac{NC^\text{t}}{NC^\text{p}_1 + kC^\text{p}_2 + C^\text{t}}$, and $C^\text{p}_1, C^\text{p}_2 \ll C^\text{t}$, we can see how cost savings improve as the number of candidate datasets $N$ enlarge. Considering that the candidate space $N$ is very large in the real-world, this framework enables practitioners to optimize pre-training datasets at a fraction of the cost.

% \paragraph{Limitation and Future Direction.} We discuss limitation and more future direction in Appendix. \ref{app:limitation}.
\subsection{Limitation and Future Direction} \label{app:limitation}

We discuss several limitations of \textsc{rBridge} that present opportunities for future research. First, frontier models do not achieve perfect accuracy on reasoning tasks, potentially resulting in imperfect extracted reasoning trace $R^\phi$. While we utilize all available $R^\phi$ without filtering and demonstrate substantial performance and efficiency gains. Our preliminary experiments showed minimal improvements from post-hoc filtering of incorrect generations. Future work could investigate more sophisticated quality assurance mechanisms, such as ensemble methods leveraging multiple frontier models.

Second, frontier models occasionally fail to produce outputs in the required format for reasoning trace $R^\phi$ extraction. Our current approach provides one additional generation attempt before excluding the sample from training. This represents a limitation that future research could address through more robust prompt engineering, or ensembling multiple frontier models.

Finally, while we present a potential framework for further practical application in \textbf{\S\ \ref{sec:discussion}}, the efficient and effective implementation of this system remains an open challenge. Future work can further explore how to best leverage the methodological contribution of \tsc{rBridge} (\textbf{\S\ \ref{sec:rbridge}}) and experimental robustness (\textbf{\S\ \ref{sec:exp}}) on more complex frameworks.

% From this, they can derive the predictive functional form between \tsc{rBridge} and the target metrics. 

% While the value of $k$ can be chosen by considering the compute budget required to achieve the desired accuracy, assuming $k=10$, $N=100$, $m=$ 7B with equivalent trained tokens across model sizes, and the mentioned proxy size $n=$ 0.09B results in $\sim$4\% of the cost incurred by the naïve approach of optimizing at target scale. 
